A completamento della fase di progettazione e implementazione fisica, in questo capitolo vengono presentate le query SQL corrispondenti alle attività fondamentali definite precedentemente nella Sezione~\ref{sec:operazioni}.

Come anticipato durante la raccolta dei requisiti, le operazioni elencate di seguito rappresentano un insieme selezionato di attività significative, atte a coprire in maniera equilibrata le diverse aree dello schema.
L'implementazione pratica di queste interrogazioni dimostra la capacità della base di dati di rispondere in modo robusto ed efficiente sia alle normali esigenze operative di una segreteria didattica, sia alle richieste di analisi statistica avanzata sulle carriere degli studenti.

\section{Implementazione delle operazioni di lettura}
\label{sec:implementazione-delle-operazioni-di-lettura}

\subsection{Frequenza degli insegnamenti nei piani di studio}
\label{ssec:frequenza-degli-insegnamenti-nei-piani-di-studio}
Questa interrogazione risponde all'esigenza di determinare l'insegnamento (o gli insegnamenti, in caso di \textit{ex aequo}) che compare più frequentemente all'interno dei piani di studio degli iscritti.

\begin{lstlisting}[language=SQL]
WITH conteggi_insegnamenti AS (
    SELECT i.nome, COUNT(*) AS frequenza
    FROM piano_di_studio pds
    JOIN insegnamento i ON pds.codice_insegnamento = i.codice
    GROUP BY i.codice, i.nome
)
SELECT nome, frequenza
FROM conteggi_insegnamenti
WHERE frequenza = (SELECT MAX(frequenza) FROM conteggi_insegnamenti);
\end{lstlisting}

L'operazione esegue inizialmente un raggruppamento per calcolare le occorrenze di ogni corso tramite una \textit{Common Table Expression} (\texttt{CTE}) denominata \texttt{conteggi\_insegnamenti}.
Successivamente, la query principale interroga questa tabella temporanea filtrando i risultati: tramite una sottoquery, vengono estratti e mostrati esclusivamente i record la cui frequenza coincide con il valore massimo assoluto.


\subsection{Estrazione del calendario settimanale di un docente}
\label{ssec:estrazione-del-calendario-settimanale-di-un-docente}
Questa operazione permette di ricostruire l'orario settimanale completo delle lezioni per un determinato professore, vincolando la ricerca a uno specifico anno accademico e al relativo periodo didattico.

\begin{lstlisting}[language=SQL]
SELECT i.nome, l.giorno, l.fascia_oraria, l.aula
FROM lezione l
JOIN insegnamento i ON l.codice_insegnamento = i.codice
JOIN insegnamento_edizione ie ON l.codice_insegnamento = ie.codice_insegnamento
    AND l.anno_accademico = ie.anno_accademico
    AND l.periodo = ie.periodo
JOIN docente d ON ie.cf_docente = d.cf
WHERE d.nome = 'Maria'
    AND d.cognome = 'Gozzi'
    AND l.anno_accademico = '2024/2025'
    AND l.periodo = '2';
\end{lstlisting}

L'interrogazione si articola su una catena di \texttt{JOIN} che collega la tabella \texttt{lezione} all'anagrafica del \texttt{docente} tramite la tabella ponte \texttt{insegnamento\_edizione}.
È rilevante notare come i vincoli di giunzione sfruttino l'intera chiave composta: il DBMS esegue i controlli temporali in modo nativo, beneficiando direttamente della presenza del periodo didattico propagato fisicamente durante la stesura dello schema relazionale.

\subsection{Calcolo del tasso di superamento degli insegnamenti}
\label{ssec:calcolo-del-tasso-di-superamento-degli-insegnamenti}
La query quantifica la percentuale di prove d'esame superate con successo rispetto al volume totale dei tentativi registrati per ciascun insegnamento.

\begin{lstlisting}[language=SQL]
WITH tentativi_totali AS (
    SELECT codice_insegnamento, COUNT(*) AS totali
    FROM esame
    GROUP BY codice_insegnamento
),
tentativi_superati AS (
    SELECT codice_insegnamento, COUNT(*) AS superati
    FROM esame
    WHERE punteggio >= 18
    GROUP BY codice_insegnamento
)
SELECT
    i.nome,
    (COALESCE(ts.superati, 0) * 100.0 / NULLIF(tt.totali, 0)) AS percentuale_superamento
FROM insegnamento i
LEFT JOIN tentativi_totali tt ON i.codice = tt.codice_insegnamento
LEFT JOIN tentativi_superati ts ON i.codice = ts.codice_insegnamento
ORDER BY percentuale_superamento DESC;
\end{lstlisting}

Per strutturare questo calcolo statistico in modo pulito e modulare, il problema logico è stato scomposto utilizzando due \textit{Common Table Expressions} (\texttt{WITH}).
La prima tabella temporanea, \texttt{tentativi\_totali}, calcola il divisore aggregando tutte le prove, mentre la seconda, \texttt{tentativi\_superati}, agisce da dividendo filtrando esclusivamente le valutazioni sufficienti ($\ge 18$).

La query principale parte dalla tabella \texttt{insegnamento} e utilizza delle \texttt{LEFT JOIN} per mantenere tutti i corsi nel risultato finale, compresi quelli per cui non è registrato alcun esame o alcuna promozione. Per gestire correttamente questi casi, la funzione \texttt{COALESCE} trasforma le assenze di superamenti in 0, mentre \texttt{NULLIF} previene il classico errore di divisione per zero nel caso in cui i tentativi totali siano assenti.

\subsection{Media dei tentativi per gli studenti promossi}
\label{ssec:media-dei-tentativi-per-gli-studenti-promossi}
Questa interrogazione analizza l'effettiva complessità didattica di un insegnamento, calcolando il numero medio di tentativi effettuati dagli studenti che sono riusciti a superare con successo la materia.

\begin{lstlisting}[language=SQL]
WITH tentativi_per_passare (matricola, num_tentativi) AS (
    SELECT e.matricola, COUNT(*)
    FROM esame e
    JOIN insegnamento i ON e.codice_insegnamento = i.codice
    WHERE i.nome = 'Teoria dei segnali'
    GROUP BY e.matricola
    HAVING MAX(e.punteggio) >= 18
)
SELECT AVG(num_tentativi) AS media_tentativi
FROM tentativi_per_passare;
\end{lstlisting}

\section{Implementazione delle operazioni di modifica}
\label{sec:implementazione-delle-operazioni-di-modifica}

Le operazioni di aggiornamento e rimozione alterano in modo permanente lo stato della base di dati.
Tali procedure richiedono l'utilizzo di blocchi transazionali protetti o di sub-query logiche per non compromettere la coerenza referenziale dello schema.

\subsection{Gestione della rinuncia agli studi}
\label{ssec:gestione-della-rinuncia-agli-studi}
L'eliminazione dei dati relativi a uno studente in seguito a una rinuncia formale agli studi richiede un'attenzione tecnica specifica e non può risolversi in una singola e semplice istruzione di rimozione.

\begin{lstlisting}[language=SQL]
BEGIN;
-- Rimozione preventiva obbligatoria a causa del vincolo RESTRICT
DELETE FROM esame
WHERE matricola = 15;

-- Cancellazione del profilo anagrafico (propaga il CASCADE su piano_di_studio)
DELETE FROM studente
WHERE matricola = 15;
COMMIT;
\end{lstlisting}

Come definito a livello di schema fisico, la tabella \texttt{esame} è rigorosamente protetta dal vincolo \texttt{ON DELETE RESTRICT} per tutelare i registri da cancellazioni accidentali. È quindi obbligatorio incapsulare i comandi all'interno di una transazione. Questa procede prima alla rimozione manuale dello storico dei voti e solo successivamente cancella l'anagrafica dello studente.
Quest'ultimo comando elimina contestualmente e in modo automatico i record associati all'interno della tabella \texttt{piano\_di\_studio}, per effetto del vincolo \texttt{ON DELETE CASCADE}.

\subsection{Sostituzione in blocco di una cattedra annuale}
\label{ssec:sostituzione-in-blocco-di-una-cattedra-annuale}
In caso di necessità organizzative, la segreteria dell'ateneo deve poter trasferire istantaneamente la titolarità delle edizioni annuali da un docente a un nuovo sostituto.

\begin{lstlisting}[language=SQL]
UPDATE insegnamento_edizione
SET cf_docente = 'ZCCNTL28S02G280O'
WHERE cf_docente = 'PTRFBA87S13G935U'
  AND anno_accademico = '2025/2026'
  AND EXISTS (
      SELECT 1
      FROM abilitazione_docente_insegnamento adi
      WHERE adi.cf_docente = 'ZCCNTL28S02G280O'
        AND adi.codice_insegnamento = insegnamento_edizione.codice_insegnamento
  );
\end{lstlisting}

Per garantire che l'operazione di \texttt{UPDATE} massivo non assegni cattedre in modo illegittimo o a docenti non qualificati, è stata inserita una rigorosa condizione di validazione tramite la sottoquery \texttt{EXISTS}.
Il DBMS applicherà le modifiche ai record nella tabella \texttt{insegnamento\_edizione} esclusivamente se il codice fiscale del nuovo docente risulta già formalmente associato all'insegnamento in questione all'interno della tabella di riferimento \texttt{abilitazione\_docente\_insegnamento}, garantendo totale aderenza alle regole dell'ateneo.